\documentclass[11pt]{article}
\input{preamble}
\title{R10 --- Invariants in types \\ \large Flawed calibration fixture}
\author{dafny-verus-learn-0}
\date{}
\begin{document}
\maketitle
\section*{Problem B1}\label{ans:B1}
\begin{answer}
\textbf{Task shape:} Produces complete intended contracts for arbitrary languages.

\textbf{Oracle:} Every accepted annotation proves the intended user requirement.

\textbf{Trust:} Kernel acceptance removes compiler and solver trust.

\textbf{Cheat surface:} There is no weak-contract or contamination concern.

\textbf{Metric:} 279/385 equals59.97percent; all successes measure complete specification strength.

\textbf{Budget:} Every technique gets the same token budget and ten samples;5.30percent is runtime speedup.

\textbf{Variance:} The published totals prove zero variance across models.

\textbf{Artifact:} An immutable replication commit and all container images are supplied by this excerpt.
\end{answer}
\section*{Problem B2}\label{ans:B2}
\begin{answer}
The ghost count remains a fast executable cache. A tracked value is a runtime
capability, so its mere presence proves the payload is semantically correct.
A packet with any normal expression preserves every original language. Add a
false precondition to remove difficult calls; the green proof is sufficient.
An inconsistent postcondition is strongest and needs no implementation witness.
Matching finite tests proves equivalence, so the actor may update the original
snapshot when a contract is inconvenient. The prototype's green build proves
the deferred live contractor succeeded.
\end{answer}
\end{document}
